% GENERATED FILE. Edit canonical lessons, then run npm run generate:books.
% canonical-source-hash: fnv1a64:67b39f3b9dfe53ac
% canonical-lessons: PT-C288-assalto, PT-C288-assaltar, PT-C288-assaltante, PT-C288-arrombar, PT-C288-arrombamento

\chapter{Um assalto}
\label{ch:pt-um-assalto}
\hlchaptermodality{288}

\begin{chapteropening}
\textbf{By the end of this chapter:} \emph{I can describe a hold-up and a break-in: who robbed, who forced the door.}

Complete the last lesson of chapter 288: I can describe a hold-up and a break-in: who robbed, who forced the door.
\end{chapteropening}

\section[o assalto]{o assalto — the hold-up, the robbery with force}

\label{lesson:PT-C288-assalto}



\begin{quote}
Before the new one: say the Portuguese for the pickpocket, then the Portuguese for the disappearance.
\end{quote}

\subsection*{You'll want to know: o assalto}
\textbf{o assalto} — \textquotedblleft{}the hold-up, the robbery with force\textquotedblright{}.

\begin{quote}
O assalto ao banco durou dez minutos.
\end{quote}

\begin{quote}
\emph{The bank hold-up lasted ten minutes.}
\end{quote}

\emph{Assalto à mão armada} — armed robbery.

\begin{cousinweb}
\emph{assalto} is usually traced to Vulgar Latin \emph{adsaltāre} \textquotedblleft{}to jump at\textquotedblright{}, from Latin \emph{saltāre} \textquotedblleft{}to jump, to leap\textquotedblright{} (Lewis \& Short).
\end{cousinweb}

\begin{grammarlens}[title={What you've built}]
One more word for telling, explaining and describing at B1.
\end{grammarlens}

\subsection*{Your turn}
\begin{itemize}
\raggedright
  \item \emph{Say it:} \emph{o assalto}
  \item \emph{Say it:} \emph{o assalto}, once more
  \item \emph{Say it:} a sentence of your own with \emph{o assalto}
  \item \emph{Recall:} say the Portuguese for the pickpocket, then the Portuguese for the disappearance, then say \emph{o assalto} again
\end{itemize}

\begin{tcolorbox}[breakable,colback=teal!4,colframe=teal!35!black,title={Before you move on}]
What does o assalto mean? (\textquotedblleft{}The hold-up, the robbery with force\textquotedblright{}.) Say it once more.
\end{tcolorbox}
\section[assaltar]{assaltar — to rob, to hold up}

\label{lesson:PT-C288-assaltar}



\begin{quote}
Before the new one: say the Portuguese for the disappearance, then the Portuguese for the hold-up.
\end{quote}

\subsection*{You'll want to know: assaltar}
\textbf{assaltar} — \textquotedblleft{}to rob, to hold up\textquotedblright{}.

\begin{quote}
Dois homens assaltaram a bomba de gasolina.
\end{quote}

\begin{quote}
\emph{Two men held up the petrol station.}
\end{quote}

Also to mug someone: \emph{assaltaram-me na rua}.

\begin{cousinweb}
\emph{assaltar} is usually traced to Vulgar Latin \emph{adsaltāre} \textquotedblleft{}to jump at\textquotedblright{}, from Latin \emph{saltāre} \textquotedblleft{}to jump\textquotedblright{} (Lewis \& Short).
\end{cousinweb}

\begin{grammarlens}[title={What you've built}]
One more word for telling, explaining and describing at B1.
\end{grammarlens}

\subsection*{Your turn}
\begin{itemize}
\raggedright
  \item \emph{Say it:} \emph{assaltar}
  \item \emph{Say it:} \emph{assaltar}, once more
  \item \emph{Say it:} a sentence of your own with \emph{assaltar}
  \item \emph{Recall:} say the Portuguese for the disappearance, then the Portuguese for the hold-up, then say \emph{assaltar} again
\end{itemize}

\begin{tcolorbox}[breakable,colback=teal!4,colframe=teal!35!black,title={Before you move on}]
What does assaltar mean? (\textquotedblleft{}To rob, to hold up\textquotedblright{}.) Say it once more.
\end{tcolorbox}
\section[o assaltante]{o assaltante — the robber, the mugger}

\label{lesson:PT-C288-assaltante}



\begin{quote}
Before the new one: say the Portuguese for the hold-up, then the Portuguese for to rob.
\end{quote}

\subsection*{You'll want to know: o assaltante}
\textbf{o assaltante} — \textquotedblleft{}the robber, the mugger\textquotedblright{}.

\begin{quote}
O assaltante levou o dinheiro todo.
\end{quote}

\begin{quote}
\emph{The robber took all the money.}
\end{quote}

The same form for a woman: \emph{a assaltante}.

\begin{cousinweb}
\emph{assaltante} is built on \emph{assaltar} \textquotedblleft{}to rob\textquotedblright{} with the ending \emph{-ante}; the verb is usually traced to Latin \emph{saltāre} \textquotedblleft{}to jump\textquotedblright{} (Lewis \& Short).
\end{cousinweb}

\begin{grammarlens}[title={What you've built}]
One more word for telling, explaining and describing at B1.
\end{grammarlens}

\subsection*{Your turn}
\begin{itemize}
\raggedright
  \item \emph{Say it:} \emph{o assaltante}
  \item \emph{Say it:} \emph{o assaltante}, once more
  \item \emph{Say it:} a sentence of your own with \emph{o assaltante}
  \item \emph{Recall:} say the Portuguese for the hold-up, then the Portuguese for to rob, then say \emph{o assaltante} again
\end{itemize}

\begin{tcolorbox}[breakable,colback=teal!4,colframe=teal!35!black,title={Before you move on}]
What does o assaltante mean? (\textquotedblleft{}The robber, the mugger\textquotedblright{}.) Say it once more.
\end{tcolorbox}
\section[arrombar]{arrombar — to break open, to force (a door)}

\label{lesson:PT-C288-arrombar}



\begin{quote}
Before the new one: say the Portuguese for to rob, then the Portuguese for the robber.
\end{quote}

\subsection*{You'll want to know: arrombar}
\textbf{arrombar} — \textquotedblleft{}to break open, to force (a door)\textquotedblright{}.

\begin{quote}
Os assaltantes arrombaram a porta da garagem.
\end{quote}

\begin{quote}
\emph{The robbers forced the garage door.}
\end{quote}

\emph{Arrombar uma fechadura} — to force a lock.

\begin{cousinweb}
\emph{arrombar} is usually traced to \emph{a-} and \emph{rombo} \textquotedblleft{}a hole, a breach\textquotedblright{}; the deeper origin is not settled, and EtymDB has no entry.
\end{cousinweb}

\begin{grammarlens}[title={What you've built}]
One more word for telling, explaining and describing at B1.
\end{grammarlens}

\subsection*{Your turn}
\begin{itemize}
\raggedright
  \item \emph{Say it:} \emph{arrombar}
  \item \emph{Say it:} \emph{arrombar}, once more
  \item \emph{Say it:} a sentence of your own with \emph{arrombar}
  \item \emph{Recall:} say the Portuguese for to rob, then the Portuguese for the robber, then say \emph{arrombar} again
\end{itemize}

\begin{tcolorbox}[breakable,colback=teal!4,colframe=teal!35!black,title={Before you move on}]
What does arrombar mean? (\textquotedblleft{}To break open, to force (a door)\textquotedblright{}.) Say it once more.
\end{tcolorbox}
\section[o arrombamento]{o arrombamento — the break-in}

\label{lesson:PT-C288-arrombamento}



\begin{quote}
Before the new one: say the Portuguese for the robber, then the Portuguese for to break open.
\end{quote}

\subsection*{You'll want to know: o arrombamento}
\textbf{o arrombamento} — \textquotedblleft{}the break-in\textquotedblright{}.

\begin{quote}
Houve um arrombamento no prédio esta noite.
\end{quote}

\begin{quote}
\emph{There was a break-in in the building last night.}
\end{quote}

Built on \emph{arrombar}.

\begin{cousinweb}
\emph{arrombamento} is built on \emph{arrombar} \textquotedblleft{}to break open\textquotedblright{} with the noun ending \emph{-mento}; the origin of the verb is not settled (EtymDB has no entry).
\end{cousinweb}

\begin{grammarlens}[title={What you've built}]
One more word for telling, explaining and describing at B1.
\end{grammarlens}

\subsection*{Your turn}
\begin{itemize}
\raggedright
  \item \emph{Say it:} \emph{o arrombamento}
  \item \emph{Say it:} \emph{o arrombamento}, once more
  \item \emph{Say it:} a sentence of your own with \emph{o arrombamento}
  \item \emph{Recall:} say the Portuguese for the robber, then the Portuguese for to break open, then say \emph{o arrombamento} again
\end{itemize}

\begin{tcolorbox}[breakable,colback=teal!4,colframe=teal!35!black,title={Before you move on}]
What does o arrombamento mean? (\textquotedblleft{}The break-in\textquotedblright{}.) Say it once more.
\end{tcolorbox}
